\documentclass[10pt,twocolumn,letterpaper]{article}

\usepackage[pagenumbers]{cvpr}
\usepackage{amsmath}
\usepackage{amssymb}
\usepackage{booktabs}
\usepackage{graphicx}
\usepackage{microtype}
\usepackage{multirow}
\usepackage{placeins}
\usepackage{tabularx}
\usepackage{xcolor}
\usepackage{url}

\definecolor{cvprblue}{rgb}{0.21,0.49,0.74}
\usepackage[colorlinks=true,allcolors=cvprblue]{hyperref}

\graphicspath{{../../results/final_report_figures/}}
\newcommand{\dataset}{iNaturalist 2021 Mini}
\newcommand{\topone}{top-1}
\newcommand{\topfive}{top-5}
\newcommand{\vect}[1]{\boldsymbol{#1}}
\DeclareMathOperator*{\argmax}{arg\,max}

\setlength{\tabcolsep}{4pt}


\title{Fine-Grained Species Recognition at Scale:\\
Traditional Baselines, Transfer Learning, Robustness, and Grad-CAM}

\author{
Yuan Cao \quad
Kaishi Yang\\
Aditya Winata \quad
Jialin Guo\\
Yuzhe Qiu\\
COMP9517 Computer Vision Group Project
}

\begin{document}
\maketitle
% Generated from frozen Stage 1--4 evidence. Do not edit by hand.
\newcommand{\BestTestTopOne}{0.7274}
\newcommand{\BestTestTopFive}{0.9026}
\newcommand{\BestTestMacroFOne}{0.723535}
\newcommand{\BestMeanCorruptedMacroFOne}{0.562260}
\newcommand{\PrimaryTraditionalTestTopOnePercent}{12.52}
\newcommand{\PrimaryTraditionalTestTopFivePercent}{26.78}
\newcommand{\PrimaryTraditionalTestMacroFOnePercent}{11.16}
\newcommand{\RandomForestTestTopOnePercent}{8.22}
\newcommand{\RandomForestTestTopFivePercent}{20.26}
\newcommand{\RandomForestTestMacroFOnePercent}{6.54}
\newcommand{\BoVWOnlyTestTopOnePercent}{4.76}
\newcommand{\BoVWOnlyTestTopFivePercent}{10.82}
\newcommand{\BoVWOnlyTestMacroFOnePercent}{4.14}
\newcommand{\BestTestTopOnePercent}{72.74}
\newcommand{\BestTestTopFivePercent}{90.26}
\newcommand{\BestTestMacroFOnePercent}{72.35}
\newcommand{\BestMeanCorruptedTopOnePercent}{56.34}
\newcommand{\BestMeanCorruptedMacroFOnePercent}{56.23}
\newcommand{\BestWorstCorruptedTopOnePercent}{18.94}
\newcommand{\BestWorstCorruptedMacroFOnePercent}{19.20}
\newcommand{\RandomForestMeanCorruptedTopOnePercent}{6.26}
\newcommand{\RandomForestMeanCorruptedMacroFOnePercent}{4.89}
\newcommand{\RandomForestWorstCorruptedTopOnePercent}{2.36}
\newcommand{\RandomForestWorstCorruptedMacroFOnePercent}{1.45}


\begin{abstract}
Fine-grained species recognition must separate visually similar taxa while
remaining reliable across pose, scale, background, and image quality.  We
study this problem on a deterministic 500-species subset of iNaturalist 2021
Mini containing 20,000 training, 5,000 internal-validation, and 5,000 held-out
test images.  A split-controlled protocol compares handcrafted colour, texture,
shape, and SIFT bag-of-visual-words pipelines with convolutional networks
trained from scratch and with ImageNet transfer.  The primary traditional
model is a 3,704-D LinearSVC combining a 2,424-D handcrafted descriptor with
1,280-D spatial SIFT-BoVW, reaching
\PrimaryTraditionalTestTopOnePercent\% held-out top-1 accuracy.  A pretrained
EfficientNet-B0 reaches \BestTestTopOnePercent\% top-1,
\BestTestTopFivePercent\% top-5, and \BestTestMacroFOnePercent\% macro F1.
The fixed 64-condition corruption study uses a distinct 2,424-D, 300-tree
handcrafted Random Forest without BoVW as its traditional comparator and
otherwise retains the frozen deep models.  Confusion analysis and Grad-CAM
expose fine-grained look-alikes as well as plausible colour and texture
shortcuts.  Together, the evidence establishes transfer learning as the
strongest approach in this setting while identifying robustness,
protocol-provenance, and explanation limitations that clean accuracy alone
conceals.
\end{abstract}

\section{Introduction}
\label{sec:introduction}

Species recognition is an unusually demanding image-classification problem.
The label space is large, distinctions may depend on a small anatomical
marking, and photographs are acquired opportunistically rather than under
controlled laboratory conditions.  In community-science collections, a
target can occupy only a few pixels, be partially occluded, or resemble a
closely related species while the background offers a tempting but unstable
cue.  The iNaturalist benchmarks capture these challenges at scale and have
therefore become important tests of visual representation learning
\cite{vanhorn2018inaturalist,vanhorn2021benchmarking}.

This work asks how far progressively more expressive computer-vision methods
can address a fixed 500-class recognition task.  We begin with interpretable
handcrafted representations---colour histograms, local binary patterns (LBP),
histograms of oriented gradients (HOG), and SIFT bag-of-visual-words
(BoVW)---paired with linear or tree-based classifiers.  We then compare
ResNet18 trained from random initialization against ImageNet-pretrained
ResNet18 and EfficientNet-B0 models.  Beyond clean-set accuracy, we test
controlled image corruptions and inspect decisions with Grad-CAM.  This broad
comparison is useful because a single accuracy number cannot reveal whether a
model is learning transferable morphology, memorising texture, or failing
under a realistic loss of image quality.

Equally important is evaluation discipline.  The official iNaturalist
validation archive is used only as a held-out test source.  Every fitted
transform, visual vocabulary, classifier, and neural-network update sees the
training partition alone.  Hyperparameters, early stopping, and ablations use
an internal validation partition.  Stage~3 records its selected configurations
before test access.  For the rewritten Stage~2 track, the present selection
manifest cryptographically freezes the complete selected state going forward,
but the available history cannot independently establish that an equivalent
lock existed before the already-recorded test evaluation.  We therefore do not
claim a retrospectively untouched Stage~2 test split.
Figure~\ref{fig:method-overview} summarises the intended protocol.

\begin{figure*}[t]
  \centering
  \includegraphics[width=0.98\textwidth]{method_overview.pdf}
  \caption{End-to-end project workflow.  Traditional and deep models share
  the same frozen class mapping and partitions.  Validation controls
  model/configuration selection and train-only pilots control feasibility;
  current artifacts are hash-bound before any future test access.}
  \label{fig:method-overview}
\end{figure*}

Our contributions are fourfold.  First, we construct a deterministic,
audited 500-species protocol with 40 training, 10 validation, and 10 test
images per species.  Second, we provide complete traditional and deep
baselines under common metrics, including training and inference timing.
Third, we compare feature components, ImageNet pretraining, augmentation
strength, and backbone freezing on validation data; the matched ResNet18
augmentation pair provides the strict one-factor ablation.
Fourth, we complement clean evaluation with corruption robustness, confusion
analysis, failure cases, and target-specific Grad-CAM visualisations.  The
result is not merely a leaderboard: it is evidence about which design choices
matter, where performance breaks down, and which apparent explanations remain
uncertain.

\section{Literature Review}
\label{sec:literature}

\paragraph{Fine-grained recognition and natural-world data.}
Fine-grained visual categorisation differs from generic object recognition
because between-class variation can be smaller than within-class variation.
Pose, life stage, sex, season, and viewpoint may transform an organism's
appearance, while sibling species differ only in local structure.  The
iNaturalist 2018 dataset established a large, naturally imbalanced benchmark
for species classification and detection \cite{vanhorn2018inaturalist}.  The
2021 benchmark expanded this setting and highlighted the importance of
representations that transfer across the long-tailed visual world
\cite{vanhorn2021benchmarking}.  Our balanced subset does not reproduce the
full long tail, but retains the taxonomic diversity and uncontrolled image
conditions that make the benchmark difficult.

\paragraph{Handcrafted representations.}
Before end-to-end representation learning, recognition pipelines separated
feature design from classification.  SIFT describes local gradient structure
and is robust to scale and rotation changes \cite{lowe2004distinctive}; BoVW
quantises local descriptors into a fixed-length image histogram
\cite{csurka2004visual}.  HOG aggregates oriented gradients over spatial cells
and captures object shape \cite{dalal2005histograms}, whereas uniform LBP
encodes local texture transitions \cite{ojala2002multiresolution}.  Colour
histograms add information these grayscale descriptors omit.  Linear support
vector machines learn a maximum-margin decision boundary
\cite{cortes1995support}, and random forests combine decorrelated decision
trees \cite{breiman2001random}.  These methods offer low-cost, interpretable
reference points, but their fixed statistics cannot adapt to the hierarchy of
parts, textures, and context needed for hundreds of species.

\paragraph{Deep transfer learning.}
ImageNet enabled supervised pretraining at a scale sufficient to learn broadly
useful visual primitives \cite{deng2009imagenet}.  Residual connections make
deeper networks easier to optimise \cite{he2016deep}, while EfficientNet uses
compound depth/width/resolution scaling to improve parameter efficiency
\cite{tan2019efficientnet}.  Features learned on a source task can transfer to
a related target, although transferability generally decreases with task and
layer distance \cite{yosinski2014transferable}.  Data augmentation reduces
overfitting by encoding nuisance invariances; random cropping, colour
perturbation, and erasing are common choices
\cite{shorten2019survey,zhong2020random}.  We test both full fine-tuning and a
frozen-backbone linear-head regime, because the latter reveals how much of the
target task is already linearly separable in ImageNet features.

\paragraph{Robustness and explanation.}
Standard test sets can conceal brittleness to blur, noise, and compression.
Image-corruption benchmarks formalise this concern by evaluating fixed models
under parameterised degradations \cite{hendrycks2019benchmarking}.  Our study
adopts the same principle, though with a project-specific corruption suite and
severity schedule applied to every held-out image.  Grad-CAM instead probes
spatial evidence: gradients of a target score weight convolutional feature
maps to produce a coarse localisation heatmap \cite{selvaraju2017gradcam}.
Such maps are useful diagnostic evidence, but they do not prove causal feature
use or explanation faithfulness.  We therefore interpret them jointly with
persisted predictions, true-target maps, and confusion statistics rather than
as standalone explanations.

\section{Dataset and Protocol}
\label{sec:dataset}

\subsection{Frozen 500-species subset}

We use \dataset{} \cite{vanhorn2021benchmarking}.  From the COCO-style
annotations, a class is eligible only if it contains at least 50 images in
\texttt{train\_mini} and 10 images in the official validation archive.  All
10,000 categories satisfy this threshold.  We sort original category IDs,
sample exactly 500 without replacement using NumPy PCG64 with seed 42, then
sort the selected IDs and map them contiguously to labels 0--499.  For each
class, a class-specific deterministic permutation selects 50
\texttt{train\_mini} images: 40 become training data and 10 become internal
validation data.  Ten images from the official validation archive form the
held-out test partition.

The resulting manifests contain 20,000 training, 5,000 validation, and 5,000
test rows.  Scientific/common names, original category IDs, image IDs, source
splits, and portable local paths are retained.  Automated validation checks
exact per-class counts, contiguous labels, annotation provenance, safe paths,
file existence, and duplicates by path, image ID, and UUID filename.  No
duplicate or cross-partition record was observed.  Figure~\ref{fig:dataset}
shows deterministic samples and illustrates the substantial changes in
object scale, viewpoint, background, and taxonomic kingdom.

\begin{figure}[t]
  \centering
  \includegraphics[width=\columnwidth]{dataset_examples.jpg}
  \caption{Examples from the frozen partitions.  Images are EXIF-oriented for
  decoding; labels and assignments remain unchanged.}
  \label{fig:dataset}
\end{figure}

\subsection{Leakage controls and test isolation}

All preprocessing that can learn from data is fitted on the training manifest
only.  This includes feature scaling, the SIFT vocabulary, linear and forest
classifiers, and neural-network optimisation.  Validation images determine
early stopping, feature ablations, and the choice between pretrained
augmentation variants.  Train-only pilots determine solver, memory, and
runtime feasibility.  The official test manifest is not read by training or
validation commands.  Stage~3 binds selected checkpoints and configurations to
cryptographic hashes before test evaluation.  The restored Stage~2 selection
gate now enforces the same ordering for all future evaluations, although the
repository cannot independently prove that this gate preceded the traditional
test artifacts already present when the rewrite was integrated.
Robustness and Grad-CAM operate only after this clean-test evaluation and never
feed changes back into a model.  This protocol prevents direct training-time
image leakage and future iteration on held-out scores; the retrospective
Stage~2 chronology limitation remains explicit.

Because the subset is class-balanced, overall accuracy and macro recall are
numerically equal here, but we retain both concepts in the evaluation code.
Balanced sampling also means every class contributes equally to macro metrics.
It does not make the task easy: each species has only 40 learning examples,
and the held-out per-class recall is quantised in increments of 0.1 because
there are ten test examples per class.

\section{Methods}
\label{sec:methods}

\subsection{Traditional computer vision}

\paragraph{Block-structured handcrafted descriptor.}
After EXIF correction and RGB conversion, each image is resized bilinearly to
$128\!\times\!128$.  Three non-negative blocks are concatenated in the fixed
order colour, texture, and HOG.  Colour uses a joint hue--saturation histogram
($12\!\times\!8$ bins) plus an eight-bin value histogram.  Texture uses
uniform LBP at $(P{=}8,R{=}1)$ and $(P{=}16,R{=}2)$.  The colour and texture
histograms are pooled over the whole image and a $2\!\times\!2$ grid, yielding
520 and 140 dimensions respectively.  HOG uses nine orientations,
$16\!\times\!16$-pixel cells, $2\!\times\!2$-cell blocks, square-root input,
and L2-Hys normalisation, producing 1,764 values.  The complete per-image
handcrafted vector is therefore 2,424-D.  Caches retain the raw blocks in
manifest order; fitted normalisation remains inside each estimator pipeline.

\paragraph{SIFT--BoVW.}
We preserve aspect ratio while limiting the longest side to 256 pixels, convert
to grayscale, and retain at most 300 SIFT keypoints per image
\cite{lowe2004distinctive}.  A 256-word visual vocabulary is learned by
mini-batch $k$-means from at most 100 descriptors from four deterministic
\emph{training} images per class.  Neither validation nor test descriptors
contribute to vocabulary fitting.  Each local descriptor uses top-3 soft
assignment to its three nearest words with Gaussian weights whose bandwidth is the
median nearest-centre distance on sampled training descriptors, reducing
visual-word boundary ambiguity \cite{vangemert2010visual}.  Counts are pooled
over the whole image and a $2\!\times\!2$ spatial grid
\cite{lazebnik2006beyond}, so the BoVW block has
$256\!\times\!5=1{,}280$ dimensions.  Images without keypoints receive an
all-zero block.

\paragraph{Normalisation and classifier selection.}
Each selected block is independently $\ell_1$ normalised, raised element-wise
to a power $\alpha$, and $\ell_2$ normalised.  At $\alpha{=}0.5$ this is the
Hellinger map \cite{arandjelovic2012three}; more generally, power
normalisation suppresses dominating histogram bins
\cite{perronnin2010improving}.  Among 19 recorded validation configurations,
the selected primary pipeline uses $\alpha{=}0.3$, block-wise normalisation,
and all four blocks.  Its 3,704-D input is the 2,424-D handcrafted descriptor
plus 1,280-D BoVW, followed by \texttt{LinearSVC} (liblinear, $C{=}1$).  The
separate BoVW-only comparator applies \texttt{LinearSVC} to 1,280 dimensions.
The Random Forest uses 300 trees, square-root feature subsampling, minimum leaf
size 2, and only the 2,424-D colour--texture--HOG descriptor.

The primary Stage~2 traditional comparison is the 3,704-D LinearSVC.  Stage~4
instead uses the 2,424-D, 300-tree handcrafted Random Forest without BoVW as
its canonical traditional robustness comparator: the corruption evaluator
recomputes per-image handcrafted descriptors and applies the fitted pipeline,
but it does not independently reconstruct the training-fitted SIFT vocabulary
required for BoVW.  The two traditional roles are therefore deliberately
distinct.

\subsection{Deep networks}

We evaluate ResNet18 \cite{he2016deep} and EfficientNet-B0
\cite{tan2019efficientnet}.  Each classifier head is replaced by a 500-output
linear layer.  The scratch ResNet18 starts from random weights.  Transfer
experiments initialise the backbone with the corresponding ImageNet weights
and randomly initialise the new head.  A frozen-head ablation holds the entire
pretrained backbone fixed and keeps its BatchNorm layers in evaluation mode;
the remaining experiments fine-tune every parameter.

At evaluation time, images are EXIF-oriented, converted to RGB, bicubically
resized to 256 pixels, centre-cropped to $224\!\times\!224$, converted to a
tensor, and normalised by ImageNet channel statistics.  Basic training uses a
random resized crop (scale 0.7--1.0) and horizontal flip.  Strong training
widens the crop scale to 0.5--1.0 and adds moderate colour jitter, a
$\pm15^{\circ}$ affine/rotation perturbation, and random erasing with
probability 0.25.  We intentionally omit vertical flips and extreme crops that
could change biologically meaningful orientation or remove diagnostic parts.

For parameters $\theta$ and a labelled mini-batch, all networks minimise
multiclass cross entropy,
\begin{equation}
  \mathcal{L}(\theta)=-\frac{1}{B}\sum_{i=1}^{B}
  \log\frac{\exp z_{i,y_i}}{\sum_{c=0}^{499}\exp z_{i,c}},
  \label{eq:crossentropy}
\end{equation}
using AdamW, cosine learning-rate annealing, weight decay $10^{-4}$, and
gradient-norm clipping at 1.0.  Effective batch size is 32.  Scratch training
uses learning rate $10^{-3}$ for at most 25 epochs; fully fine-tuned pretrained
models use $10^{-4}$ for at most 15 epochs; the frozen-head run uses
$10^{-3}$ for at most 10 epochs.  Validation macro F1 controls checkpointing,
with top-1, top-5, loss, and earlier epoch as deterministic tie-breakers.

\subsection{Robustness analysis}

Let $m_0$ be a model's clean metric and $m_{k,s}$ its metric for corruption
type $k$ at severity $s$.  We report absolute drop
$\Delta_{k,s}=m_0-m_{k,s}$, relative drop $\Delta_{k,s}/m_0$, mean corrupted
performance across 16 conditions, and the worst observed condition.
The fixed suite contains Gaussian blur, Gaussian noise, JPEG compression, and
motion blur, each at four increasing severities.  Corruptions are deterministic
and applied to all 5,000 test images; model parameters and normalisation remain
unchanged.  This isolates sensitivity to acquisition quality from retraining
or test-time adaptation.

\subsection{Grad-CAM}

For target class $c$, Grad-CAM weights convolutional channel $k$ by the
spatially averaged target-score gradient,
\begin{equation}
  \alpha_k^c=\frac{1}{Z}\sum_{u,v}
  \frac{\partial z_c}{\partial A_{uv}^{k}}, \qquad
  L^c=\operatorname{ReLU}\!\left(\sum_k\alpha_k^c A^k\right).
  \label{eq:gradcam}
\end{equation}
We compute maps at EfficientNet-B0 layer \texttt{features.8}, upsample the
activation map to the model crop, and align it back to the full image using the
recorded centre-crop geometry.  The analysis set is selected deterministically
from persisted predictions: 12 high-confidence correct cases, 12
high-confidence errors, and nine representatives of top bidirectional
confusion pairs (33 images, 54 target maps).  Correct cases use the predicted
target; errors include both predicted- and true-class targets.  This
target-specific pairing helps distinguish ``where the predicted class scores''
from ``where the true class would score'', but remains a diagnostic rather than
a causal proof.

\section{Experimental Setup}
\label{sec:setup}

\paragraph{Implementation.}
Experiments run locally on an Apple M4 Pro.  Deep models use PyTorch and
torchvision \cite{paszke2019pytorch} with MPS in FP32; traditional pipelines
use scikit-learn \cite{pedregosa2011scikit} and OpenCV
\cite{bradski2000opencv}.  Data, image, configuration, progress, plotting, and
serialisation utilities use NumPy, pandas, tqdm, scikit-image, Pillow,
Matplotlib, joblib, torchvision, and PyYAML
\cite{harris2020array,mckinney2010data,costaluis2019tqdm,vanderwalt2014scikitimage,
clark2026pillow,hunter2007matplotlib,joblib2025,marcel2010torchvision,pyyaml2025}.
Random seeds for Python, NumPy, data samplers,
workers, and model initialisation are fixed to 42.  Data decoding applies EXIF
orientation and handles grayscale or alpha-channel inputs by safe RGB
conversion.  Failed or unreadable images are fatal rather than silently
skipped.

Scratch ResNet18 uses a 25-epoch cap and patience 6.  The two pretrained
ResNet18 augmentation variants share initial pretrained weights, head
initialisation, sampler order, optimiser, scheduler, 15-epoch cap, and
patience 4; only the training transform changes.  EfficientNet-B0 uses the
same strong-augmentation optimisation policy.  A disposable preflight chooses
a feasible micro-batch without contributing weights or measurements to the
canonical run.  Before epoch one, the model is rebuilt and all random sources
are reseeded.

\paragraph{Metrics.}
For each of 500 labels we compute precision, recall, and F1 with zero assigned
when a denominator is empty, then report their unweighted macro averages.
Top-$k$ accuracy counts an example as correct when its true class is among the
$k$ largest scores.  Score columns are aligned through the estimator's
explicit class order, and stable descending ranking breaks exact ties in favour
of the lower label.  Numeric confusion matrices use rows for true labels and
columns for predictions.  We retain both overall/top-1 accuracy names for
rubric compatibility, but they are the same statistic.

Traditional training time surrounds classifier \texttt{fit} only, and its
test time surrounds scoring an already-loaded cached feature matrix.  Deep
training time includes loading, augmentation, transfer, forward/backward, and
optimiser work in the training loop, excluding validation and serialisation.
Deep test time is synchronised model-forward time; full evaluation wall time
is also logged.  Consequently, traditional and deep timing numbers describe
different pipeline boundaries and should not be interpreted as a direct
hardware-speed contest.

\paragraph{Selection and reproducibility.}
For Stage~3, validation metrics are independently recomputed from each best
checkpoint before freeze.  Test selection includes scratch ResNet18, the
validation-best pretrained ResNet18 augmentation, and completed
EfficientNet-B0; the basic-augmentation and frozen-head ResNet18 runs remain
validation-only.  For rewritten Stage~2, the selection manifest binds the 19
validation configurations, selected component order, normalisation, model and
cache hashes, train-fitted BoVW vocabulary, and Stage~1 split hashes.  Test
evaluation rejects a missing or mismatched manifest, an unlisted model, and
silent replacement of canonical metrics.  This freezes the current state
going forward.  Because no independently verifiable pre-test manifest survives
for the already-recorded traditional results, it does not establish historical
chronology.  Atomic caches, checkpoints, predictions, metric files, and
configuration hashes nevertheless make the evidence currently reported
traceable to its bound source.
Figure~\ref{fig:convergence} shows observed optimisation curves reconstructed
from the persisted epoch logs, without rerunning inference.

\begin{figure*}[t]
  \centering
  \includegraphics[width=0.92\textwidth]{training_convergence.png}
  \caption{Training loss and validation macro F1 for scratch ResNet18 and the
  best pretrained EfficientNet-B0.  Transfer learning converges faster and to
  a substantially stronger validation solution.}
  \label{fig:convergence}
\end{figure*}

\section{Experimental Results}
\label{sec:results}

\subsection{Clean validation and held-out test}

Table~\ref{tab:main-results} reports all completed baselines.  The three
traditional rows now represent distinct inputs.  The primary 3,704-D
colour--texture--HOG--BoVW LinearSVC reaches
\PrimaryTraditionalTestTopOnePercent\% test top-1,
\PrimaryTraditionalTestTopFivePercent\% top-5, and
\PrimaryTraditionalTestMacroFOnePercent\% macro F1.  The 2,424-D, 300-tree
handcrafted-only Random Forest without BoVW reaches
\RandomForestTestTopOnePercent\%,
\RandomForestTestTopFivePercent\%, and
\RandomForestTestMacroFOnePercent\%, while the 1,280-D BoVW-only LinearSVC
reaches \BoVWOnlyTestTopOnePercent\%, \BoVWOnlyTestTopFivePercent\%, and
\BoVWOnlyTestMacroFOnePercent\%.  These values are smaller than the learned
representations: scratch ResNet18 reaches 29.88\% top-1, pretrained ResNet18
reaches 67.60\%, and pretrained EfficientNet-B0 reaches
\BestTestTopOnePercent\% top-1, \BestTestTopFivePercent\% top-5, and
\BestTestMacroFOnePercent\% macro F1.  Validation/test rank agreement is
descriptive only; it is not evidence that validation overfitting or prior test
access did not occur.

% Generated from frozen evidence. Do not edit by hand.
\begin{table*}[t]
\centering
\caption{Validation and held-out test results. All values are fractions and are derived from the current hash-bound evidence. The rewritten Stage~2 configuration is frozen going forward; its pre-test locking chronology cannot be independently established retrospectively.}
\label{tab:main-results}
\small
\begin{tabularx}{\textwidth}{@{}Xccccccc@{}}
\toprule
Method & Pretr. & Val top-1 & Val top-5 & Val F1 & Test top-1 & Test top-5 & Test F1 \\
\midrule
Handcrafted + BoVW (LinearSVC) & No & 0.1294 & 0.2702 & 0.115401 & 0.1252 & 0.2678 & 0.111589 \\
Random Forest & No & 0.0860 & 0.2144 & 0.068372 & 0.0822 & 0.2026 & 0.065407 \\
SIFT-BoVW (LinearSVC) & No & 0.0472 & 0.1130 & 0.041271 & 0.0476 & 0.1082 & 0.041405 \\
ResNet18 (scratch) & No & 0.2874 & 0.5274 & 0.279552 & 0.2988 & 0.5364 & 0.292782 \\
ResNet18 (ImageNet pretrained) & Yes & 0.6754 & 0.8646 & 0.673652 & 0.6760 & 0.8624 & 0.672752 \\
EfficientNet-B0 (ImageNet pretrained) & Yes & 0.7368 & 0.9020 & 0.732628 & 0.7274 & 0.9026 & 0.723535 \\
\bottomrule
\end{tabularx}
\end{table*}

% Generated from frozen evidence. Do not edit by hand.
\begin{table*}[t]
\centering
\caption{Held-out macro metrics and measured fit/test timings. Overall accuracy equals top-1 for this single-label task; timing boundaries differ between cached-feature and end-to-end neural evaluation (Section~\ref{sec:setup}).}
\label{tab:test-details}
\small
\begin{tabularx}{\textwidth}{@{}Xccccr@{}}
\toprule
Method & Macro precision & Macro recall & Macro F1 & Train (s) & Test (s) \\
\midrule
Handcrafted + BoVW (LinearSVC) & 0.109007 & 0.125200 & 0.111589 & 1323.2 & 0.190 \\
Random Forest & 0.068343 & 0.082200 & 0.065407 & 870.5 & 1.370 \\
SIFT-BoVW (LinearSVC) & 0.038986 & 0.047600 & 0.041405 & 378.8 & 0.075 \\
ResNet18 (scratch) & 0.305342 & 0.298800 & 0.292782 & 2772.5 & 7.027 \\
ResNet18 (ImageNet pretrained) & 0.691092 & 0.676000 & 0.672752 & 1453.2 & 6.856 \\
EfficientNet-B0 (ImageNet pretrained) & 0.739866 & 0.727400 & 0.723535 & 2421.3 & 8.832 \\
\bottomrule
\end{tabularx}
\end{table*}


\begin{figure*}[t]
  \centering
  \includegraphics[width=0.92\textwidth]{final_model_comparison.png}
  \caption{Held-out comparison of the six complete methods.  The large gap
  between fixed descriptors and fine-tuned pretrained networks dominates the
  smaller within-family differences.}
  \label{fig:model-comparison}
\end{figure*}

The best network is also parameter-efficient: EfficientNet-B0 has 4.65M
trainable parameters, versus 11.43M for ResNet18, yet improves test top-1 by
5.14 points over the selected pretrained ResNet18.  Its canonical training
loop takes 2,421.3 seconds, compared with 1,453.2 seconds for that ResNet18 and
2,772.5 seconds for scratch ResNet18.  These times reflect different epoch
counts and convergence behaviour, not architecture throughput in isolation.
Figure~\ref{fig:model-comparison} visualises the clean-score separation.

\subsection{Validation-only ablations}

Table~\ref{tab:ablations} contains only validation results, so none of these
comparisons consumes additional test information.  The traditional rows show
the selected full four-block validation reference beside colour-only,
texture-only, and HOG-only controls.  The complete 19-configuration grid also
varied the power exponent, block-wise versus global normalisation, feature
subsets, and the three canonical model roles.  It selected block-wise
$\alpha{=}0.3$ with all four blocks; single-block scores must not be interpreted
as a block's contribution in the combined model.

% Generated from frozen validation evidence. Do not edit by hand.
\begin{table*}[t]
\centering
\caption{Validation-only traditional and deep-learning ablations.}
\label{tab:ablations}
\small
\begin{tabularx}{\textwidth}{@{}lXccc@{}}
\toprule
Study membership & Experiment & Val top-1 & Val top-5 & Val macro-F1 \\
\midrule
Feature reference & All blocks + BoVW & 0.1294 & 0.2702 & 0.115401 \\
Feature ablation & Colour only & 0.0608 & 0.1542 & 0.040555 \\
Feature ablation & Texture (LBP) only & 0.0356 & 0.1188 & 0.017688 \\
Feature ablation & HOG only & 0.0292 & 0.0852 & 0.019803 \\
Pretraining & ResNet18 scratch / basic & 0.2874 & 0.5274 & 0.279552 \\
Pretrain / aug. / tuning ref. & ResNet18 pretrained / basic & 0.6616 & 0.8538 & 0.660844 \\
Augmentation & ResNet18 pretrained / strong & 0.6754 & 0.8646 & 0.673652 \\
Frozen head & ResNet18 frozen backbone & 0.5746 & 0.7934 & 0.572458 \\
\bottomrule
\end{tabularx}
\end{table*}


The matched ResNet18 augmentation ablation is the cleanest deep comparison:
strong augmentation improves validation top-1 from 66.16\% to 67.54\%
(+1.38 points) and macro F1 from 66.08\% to 67.37\% (+1.28 points).  The
scratch/basic run reaches 28.74\% and the pretrained/basic run 66.16\%; this
large descriptive gap is consistent with the value of ImageNet transfer, but
is not a one-factor causal estimate because the canonical learning rate,
epoch cap, and stopping policy also differ.  Similarly, the frozen-head
configuration reaches 57.46\% versus 66.16\% for full fine-tuning, but uses a
different optimisation schedule.  It nevertheless shows that the pretrained
representation is already useful before backbone adaptation.

\subsection{Confusions and class-level failures}

The normalised 500-by-500 confusion overview in
Figure~\ref{fig:confusion} is strongly diagonal, but local off-diagonal errors
remain.  The most frequent symmetric pair is Indian Pond Heron versus Squacco
Heron, with six errors (two in one direction and four in the other).  Several
other leading pairs are taxonomically or visually close: Red-eyed Damselfly
versus Common Bluetail, Yellow-faced versus Singing Honeyeater, and
MacGillivray's versus Orange-crowned Warbler.  Figure~\ref{fig:confused-pairs}
preserves both error directions rather than collapsing them into a single
uninterpretable count.

\begin{figure}[t]
  \centering
  \includegraphics[width=0.92\columnwidth]{efficientnet_confusion.png}
  \caption{Row-normalised held-out confusion matrix for EfficientNet-B0
  (rows true, columns predicted).  The overview is complemented by the exact
  integer confused-pair counts in Figure~\ref{fig:confused-pairs}.}
  \label{fig:confusion}
\end{figure}

\begin{figure*}[t]
  \centering
  \includegraphics[width=0.88\textwidth]{top_confused_pairs.png}
  \caption{Ten most frequent bidirectional species confusions for
  EfficientNet-B0.  Colour separates the two directions of each pair.}
  \label{fig:confused-pairs}
\end{figure*}

Class-level recall ranges from 0.1 for totara, Small-leaved Blinks, and New
Zealand jasmine to 1.0 for multiple visually distinctive moth and beetle
classes.  These extremes should be interpreted cautiously because each recall
uses only ten test images.  High-confidence failures are more revealing.  For
example, a matrimony-vine image is predicted as Jerusalem cherry with 0.986
probability, and a Red Junglefowl is predicted as white pasqueflower with
0.932 probability.  Such errors motivate the spatial analysis below.

\subsection{Robustness to degradation}

Table~\ref{tab:robustness} and Figure~\ref{fig:robustness} are generated only
from a complete set of 16 corrupted conditions per model.  EfficientNet-B0 averages
\BestMeanCorruptedTopOnePercent\% top-1 and
\BestMeanCorruptedMacroFOnePercent\% macro F1; its lowest recorded top-1 is
\BestWorstCorruptedTopOnePercent\%, and its lowest recorded macro F1 is
\BestWorstCorruptedMacroFOnePercent\%.  These minima are computed independently
and need not occur under the same corruption.  The traditional curve is not
the primary 3,704-D LinearSVC.  It is the canonical 2,424-D, 300-tree
handcrafted-only Random Forest comparator without BoVW, averaging
\RandomForestMeanCorruptedTopOnePercent\%
top-1 and \RandomForestMeanCorruptedMacroFOnePercent\% macro F1, with worst
recorded top-1 and macro-F1 values of
\RandomForestWorstCorruptedTopOnePercent\% and
\RandomForestWorstCorruptedMacroFOnePercent\%, respectively; these minima are
also computed independently.  The build emits every value in
this paragraph, the table, and the curves from the same complete 64-cell
evidence matrix after the Stage~2 model change.

% Generated from frozen Stage 4 evidence. Do not edit by hand.
\begin{table*}[t]
\centering
\caption{Fixed-protocol robustness summary over 16 corrupted test conditions per model. The traditional comparator is the 2,424-D handcrafted Random Forest, not the primary 3,704-D LinearSVC.}
\label{tab:robustness}
\small
\begin{tabularx}{\textwidth}{@{}Xcccccc@{}}
\toprule
Model & Clean F1 & Mean corrupt. F1 & Worst F1 & Mean $\Delta$ top-1 & Max $\Delta$ top-1 & Rank \\
\midrule
EfficientNet-B0 pretrained & 0.723535 & 0.562260 & 0.191962 & 0.164025 & 0.538000 & 1 \\
ResNet18 pretrained & 0.672752 & 0.478144 & 0.091338 & 0.198675 & 0.587800 & 2 \\
ResNet18 scratch & 0.292782 & 0.216528 & 0.077121 & 0.076512 & 0.211200 & 3 \\
RF (2,424-D handcrafted) & 0.065407 & 0.048932 & 0.014465 & 0.019550 & 0.058600 & 4 \\
\bottomrule
\end{tabularx}
\end{table*}


\begin{figure*}[t]
  \centering
  \includegraphics[width=0.97\textwidth]{robustness_curves.png}
  \caption{Top-1 and macro-F1 trajectories under four deterministic
  corruptions.  Severity zero is the same frozen clean test evaluation for
  every curve.  The traditional curve is the 2,424-D, 300-tree handcrafted
  Random Forest without BoVW, not the primary 3,704-D LinearSVC.  No model is
  retrained or adapted.}
  \label{fig:robustness}
\end{figure*}

The absolute results demonstrate that clean accuracy is not a robustness
guarantee.  The observed pretrained deep models retain higher clean and
corrupted performance in the regenerated evidence, yet the gap between clean
and worst-condition scores remains substantial.  Robustness therefore remains
a first-order failure mode rather than a minor correction to the clean ranking.

\subsection{Grad-CAM and qualitative failure analysis}

Figure~\ref{fig:gradcam} presents predicted-target overlays drawn only from the
pre-registered analysis set.  Correct predictions often concentrate on the
organism: fruit clusters, flower heads, wing markings, or bird bodies.  This
localisation is encouraging but not sufficient evidence of the right
fine-grained cue.  The error cases are more diagnostic.  Both matrimony vine
and Jerusalem cherry maps emphasise red fruit, suggesting that shared colour
and shape dominate subtler botanical distinctions.  The Red Junglefowl map
focuses on a bright textured region and supports a plausible texture shortcut
for the flower prediction.  For the shield-bug/tick error, the predicted map
is tightly centred on the arthropod while the true-target map shifts, showing
that correct localisation can coexist with incorrect classification.

\begin{figure*}[t]
  \centering
  \includegraphics[width=0.32\textwidth]{gradcam_correct_examples.jpg}\hfill
  \includegraphics[width=0.32\textwidth]{gradcam_incorrect_examples.jpg}\hfill
  \includegraphics[width=0.32\textwidth]{gradcam_confused_pairs.jpg}
  \caption{Persisted EfficientNet-B0 Grad-CAM overlays.  Left: four
  high-confidence correct predictions.  Centre: four high-confidence errors.
  Right: representatives of frequent confused pairs.  Captions show true
  $\rightarrow$ predicted common names and softmax
  confidence.  These are qualitative diagnostics, not causal explanations.}
  \label{fig:gradcam}
\end{figure*}

Confused-pair maps likewise centre on the relevant animals, including the two
herons and two damselflies.  Their failure is therefore not simply background
attention; it is consistent with insufficient resolution or feature
specificity for closely related morphology.  Conversely, some overlays span
large rectangular crop regions, reflecting the coarse spatial resolution of
the final convolutional layer.  We consequently treat the maps as spatial
associations rather than evidence of pixel-level causation.

\section{Discussion}
\label{sec:discussion}

The clearest result is the importance of learned, transferred
representations.  Fixed colour/texture/gradient summaries provide a valid
baseline but lose spatial and semantic structure; SIFT-BoVW adds local
invariance without closing this gap.  Scratch ResNet18 learns far stronger
features from 40 examples per class, yet pretrained ResNet18 more than doubles
its top-1 accuracy.  EfficientNet-B0 then improves accuracy with fewer
parameters than ResNet18, suggesting that its architecture and pretrained
representation are well matched to the available compute and data.  Because
the scratch/pretrained schedules differ, this evidence is consistent with a
pretraining benefit but does not isolate that benefit causally.

The traditional rewrite also sharpens what the classical comparison means.
Joint colour statistics, multi-scale texture, HOG, and spatial soft-assignment
BoVW are complementary enough that the 3,704-D LinearSVC is the primary
Stage~2 traditional model.  The Stage~4 Random Forest is a different,
2,424-D, 300-tree handcrafted-only comparator without BoVW.  It remains useful
for a fixed corruption study because that evaluator can deterministically
recompute its per-image input and reuse the fitted normaliser, whereas it does
not reconstruct the training-fitted visual vocabulary.  Its robustness curve
must therefore not be read as the robustness of the best-scoring traditional
model.

The strict augmentation comparison is more modest but actionable.  Strong
augmentation improves the pretrained ResNet18 validation score without
altering its initial weights, sampling order, optimiser, or stopping policy.
The result is consistent with nuisance variability in iNaturalist images and
supports the final use of colour, affine, crop, and erasing perturbations.
However, the clean improvement does not make the model invariant to severe
degradation.  Motion blur removes high-frequency edges and local markings that
may be diagnostic, plausibly contributing to its large top-1 fall.  Training
with blur or evaluating corruption-aware
augmentation would be natural future work, but doing so after observing the
held-out corruption scores would require a new validation protocol and test
set.

Confusion and Grad-CAM analyses reveal two distinct error regimes.  Some
fine-grained pairs are correctly localised yet remain inseparable, indicating
that global $224\!\times\!224$ classification may discard diagnostic detail.
Other high-confidence failures align with conspicuous colour or texture, which
may be predictive in training but not species-specific.  A promising extension
would combine higher-resolution or part-based attention with taxonomic
hierarchy, while explicitly testing whether explanations remain stable under
background and corruption interventions.

Several limitations bound our conclusions.  The balanced 500-class subset is
reproducible but does not measure full 10,000-class or long-tail performance.
Forty training images per class constrain scratch learning, and ten test images
per class make per-class metrics noisy.  We evaluate one seed and a limited
hyperparameter budget on one M4 Pro, so small model differences may not be
statistically stable.  Traditional cached-feature timing and deep forward
timing have different scopes.  The corruption suite covers four synthetic
degradations, not distribution shifts in geography, season, camera, or
observer behaviour.  Finally, Grad-CAM is coarse and gradient-based; its maps
are useful for forming hypotheses, not certifying faithful reasoning.  A
further provenance limitation is retrospective: the integrated repository
contains validation and test evidence for rewritten Stage~2, but no surviving
artifact independently proves that the selected configuration was immutably
locked before the first test access.  The restored selection manifest freezes
the current hashes, allowed models, and validation grid going forward; it does
not manufacture a historical chronology.

Despite these limitations, the project has two practical strengths.  First,
the current split, model, prediction, metric, and selection artifacts are
cross-hashed, so subsequent drift or unauthorised test access is detectable.
Second,
the combination of clean metrics, macro scores, classwise errors, corruptions,
and explanations gives a more complete account than accuracy alone.  The
best model is genuinely useful on clean images, but the robustness and failure
analyses show exactly why it should not yet be treated as a dependable field
identifier without uncertainty handling and expert review.

\section{Conclusion}
\label{sec:conclusion}

We developed and evaluated a complete 500-species recognition pipeline on a
deterministic iNaturalist subset.  The rewritten primary traditional model
combines colour, texture, HOG, and spatial SIFT-BoVW in a 3,704-D LinearSVC and
reaches \PrimaryTraditionalTestTopOnePercent\% held-out top-1.  ResNet18
trained from scratch reaches 29.88\%, while ImageNet transfer raises the best
ResNet18 result to 67.60\%.  Pretrained EfficientNet-B0 is the strongest model
at \BestTestTopOnePercent\% top-1, \BestTestTopFivePercent\% top-5, and
\BestTestMacroFOnePercent\% macro F1.  A matched ablation finds a 1.38-point
validation gain from stronger augmentation.

The advanced studies qualify that success.  The corruption analysis uses a
distinct 2,424-D, 300-tree handcrafted Random Forest without BoVW as the
traditional comparator and derives every aggregate only after all 64 cells are
regenerated.  Confusion pairs concentrate among visually similar species, and
Grad-CAM examples suggest both meaningful organism localisation and possible
colour/texture shortcuts.  The main lesson is therefore two-sided: transfer
learning and efficient convolutional architecture provide a large practical
advance, but reliability requires evaluation beyond the clean test set.  The
current selection manifest, persisted predictions, and reproducible analyses
provide a sound basis for future higher-resolution, corruption-aware, and
taxonomy-informed methods, while the retrospective Stage~2 locking limitation
remains part of the evidence record.


\FloatBarrier
\label{LastPageMainText}
\clearpage
{\small
\bibliographystyle{ieeenat_fullname}
\bibliography{references}
}

\end{document}
